\documentclass[10pt,a4paper]{amsart}
\usepackage[margin=19mm]{geometry}
\usepackage{amsmath,amssymb,mathtools}
\usepackage[scr=rsfs,cal=euler]{mathalfa}
\usepackage{xfrac}
\usepackage[unicode,hidelinks,bookmarks=false]{hyperref}
\newcommand{\ie}{i.e.\ }
\newcommand{\cf}{cf.\ }
\newcommand{\resp}{respectively\ }
\newcommand{\Subsection}{\subsection}
\providecommand{\nobreakdash}{\nobreak\hbox{-}}
\setlength{\emergencystretch}{2em}
\multlinegap0pt
\usepackage{bm}
\usepackage{bbm}
\usepackage{dsfont}
\usepackage{xspace}
\usepackage{cancel}
\usepackage{mathtools}
\usepackage{cleveref}
\usepackage{amsthm}
\makeatletter
\@ifundefined{theorem}{\newtheorem{theorem}{Theorem}}{}
\@ifundefined{corollary}{\newtheorem{corollary}[theorem]{Corollary}}{}
\@ifundefined{lemma}{\newtheorem{lemma}[theorem]{Lemma}}{}
\makeatother
\makeatletter
\newcommand{\Fix}{\ensuremath{\operatorname{Fix}}}
\newcommand{\Id}{\ensuremath{\operatorname{Id}}}
\newcommand{\bz}{\ensuremath{\mathbf{z}}}
\expandafter\ifx\csname deflist\endcsname\relax
\newenvironment{deflist}[1][\quad]%
{\begin{list}{}{\renewcommand{\makelabel}[1]{\textrm{##1~}\hfil}%
\settowidth{\labelwidth}{\textrm{#1~}}%
\setlength{\leftmargin}{\labelwidth+\labelsep}}}%requires macro calc.sty
\fi
\newcommand{\kkk}{\ensuremath{{k\in{\mathbb N}}}}
\newcommand{\menge}[2]{\big\{{#1}~\big |~{#2}\big\}}
\newcommand{\ran}{\ensuremath{\operatorname{ran}}}
\newcommand{\tmt}{\ensuremath{T_{\text{\scriptsize MT}}}}
\makeatother
\title[The splitting algorithms by Ryu, by Malitsky-Tam, and by Cam]{The splitting algorithms by Ryu, by Malitsky-Tam, and by Campoy applied to normal cones of linear subspaces converge strongly to the projection onto the intersection}
\author{Heinz H. Bauschke, Shambhavi Singh, Xianfu Wang}
\date{Source: arXiv:2203.03832v2 (CC BY 4.0)}
\begin{document}
\maketitle

\noindent\emph{Source and licence.} The splitting algorithms by Ryu, by Malitsky-Tam, and by Campoy applied to normal cones of linear subspaces converge strongly to the projection onto the intersection. Version arXiv:2203.03832v2 (CC BY 4.0), distributed under the Creative Commons Attribution 4.0 International licence, as recorded in the arXiv licence metadata for this exact version and shown on the abstract page https://arxiv.org/abs/2203.03832v2. Full rights notice: licenses/arxiv-2203.03832-CC-BY-4.0.txt.\par\smallskip

\noindent\textit{Editorial scope.} Counted unit: the Theorem whose source label is \texttt{e:210817h} in the article (source0.tex lines 1613--1646, source label \texttt{e:210817h}), delivered together with the article's own complete original proof of it (source0.tex lines 1647--1672, one \texttt{proof} environment of 26 lines). No mathematics is written, repaired or re-derived: statement and proof are the article's own text line for line. Required context retained uncounted: c:key (lines 767--772); e:210817a (lines 1035--1039); e:linMTM (lines 1417--1445); e:linMT (lines 1447--1458); mtfixlemma (lines 1462--1507); e:210817g (lines 1502--1505). Every definition, display and label of those lines is retained unchanged. Omitted from this package and not redistributed: 1--766, 773--1034, 1040--1416, 1446--1446, 1459--1461, 1506--1612, 1673--3096. Nothing inside a retained excerpt is omitted or altered: every retained body is delivered exactly as the article's lines are written, so no omission receipt of any kind accompanies this package. Citations: the retained excerpts carry no citation command at all, so no bibliography entry of the article is transcribed and no reference is redistributed. Reference adaptation: none was needed and none was made; every reference inside the delivered unit resolves inside it, so no unresolved reference or citation is produced. Printed slips kept literal: no printed word, symbol, hypothesis or number of the counted statement or of its proof was altered. Layout and class: the article's own class is \texttt{article} with packages including float, graphicx, nameref, enumitem, mathtools, empheq, optional, xcolor, hyperref, soul, mathpazo, stmaryrd, amssymb, amsthm; the wrapper is the lane's pinned \texttt{amsart} editorial wrapper at 10pt on A4, which restates the theorem-like environments the retained text needs and adds the article's own macro definitions that the retained text needs (\texttt{\textbackslash Fix}, \texttt{\textbackslash Id}, \texttt{\textbackslash bz}, \texttt{\textbackslash kkk}, \texttt{\textbackslash menge}, \texttt{\textbackslash ran}, \texttt{\textbackslash tmt}), with extra wrapper packages (bm, bbm, dsfont, xspace, cancel, mathtools, cleveref). The delivered numbering is the unit's own. Assets: no retained span contains a figure environment, a graphics-inclusion command, a PDF-page inclusion, a table or a TikZ picture; no image file is copied into this package, the package declares no asset dependency and no image file is redistributed with it. Licence: version arXiv:2203.03832v2 of this article is distributed under the Creative Commons Attribution 4.0 International licence, as recorded in the arXiv licence metadata for this exact version and shown on the abstract page https://arxiv.org/abs/2203.03832v2. A full rights notice is delivered at licenses/arxiv-2203.03832-CC-BY-4.0.txt. Characters: every retained excerpt is pure ASCII, so the delivered engine, XeLaTeX with its default Latin Modern text font, reports no missing character.
\begin{corollary}
\label{c:key}
Let $L\colon X\to X$ be linear and averaged nonexpansive.
Then 
$(\forall x\in X)$ $L^kx \to P_{\Fix L}(x)$.
\end{corollary}

\begin{equation}
\label{e:210817a}
(\forall (i,j)\in\{1,\ldots,n\}^2)\quad 
(Q_i-Q_j)M\bz_k\to 0. 
\end{equation}

\begin{subequations}
\label{e:linMTM}
\begin{align}
M\colon X^{n-1}
&\to 
X^n\colon
\begin{pmatrix}
z_1\\
\vdots\\
z_{n-1}
\end{pmatrix}
\mapsto
\begin{pmatrix}
x_1\\
\vdots\\
x_{n-1}\\
x_n
\end{pmatrix},
\quad\text{where}\;\;
\\[+2mm]
&(\forall i\in\{1,\ldots,n\})
\;\;
x_i = \begin{cases}
P_{1}(z_1), &\text{if $i=1$;}\\
P_{i}(x_{i-1}+z_i-z_{i-1}), &\text{if $2\leq i\leq n-1$;}\\
P_{n}(x_1+x_{n-1}-z_{n -1}), &\text{if $i=n$}
\end{cases}
\end{align}
\end{subequations}

\begin{equation}
\label{e:linMT}
T := \tmt \colon X^{n-1}\to X^{n-1}\colon 
\bz\mapsto 
\bz+ 
\begin{pmatrix}
(Q_2-Q_1)M\bz\\
(Q_3-Q_2)M\bz\\
\vdots\\
(Q_n-Q_{n-1})M\bz
\end{pmatrix}. 
\end{equation}

\begin{lemma}
\label{mtfixlemma}
The fixed point set of the MT operator $T=\tmt$ is 
\begin{align}
\label{e:210817e}
\Fix T &= \menge{(z,\ldots,z)\in X^{n-1}}{z\in Z} \oplus E,
\end{align}
where 
\begin{subequations}
\label{e:bloodyE}
\begin{align}
E &:= 
\ran\Psi \cap \big(X^{n-2}\times U_n^\perp)\\
&\subseteq 
U_1^\perp \times 
\cdots
\times (U_1^\perp+\cdots+U_{n-2}^\perp)
\times \big((U_1^\perp+\cdots+U_{n-1}^\perp)\cap U_n^\perp\big)
\end{align}
\end{subequations}
and 
\begin{subequations}
\label{e:bloodyPsi}
\begin{align}
\Psi\colon U_1^\perp \times \cdots \times U_{n-1}^\perp
&\to X^{n-1}\\
(y_1,\ldots,y_{n-1})&\mapsto 
(y_1,y_1+y_2,\ldots,y_1+y_2+\cdots+y_{n-1})
\end{align}
\end{subequations}
 is the continuous linear partial sum operator which has closed range. 

Let $\bz = (z_1,\ldots,z_{n-1})\in X^{n-1}$,
and set $\bar{z} := (z_1+z_2+\cdots+z_{n-1})/(n-1)$.
Then
\begin{equation}
\label{e:210817d}
P_{\Fix T}\bz = (P_Z\bar{z},\ldots,P_Z\bar{z}) \oplus P_E\bz \in X^{n-1}
\end{equation}
and hence 
\begin{equation}
\label{e:210817g}
P_1(Q_1P_{\Fix T})\bz = P_Z\bar{z}.
\end{equation}

\end{lemma}

\begin{theorem} {\bf (main result on Malitsky-Tam splitting)} 
Given $0<\lambda<1$ and $\bz_0=(z_{0,1},\ldots,z_{0,n-1}) \in X^{n-1}$, 
generate the sequence $(\bz_k)_\kkk$ via\footnote{Recall \cref{e:linMTM} 
and \cref{e:linMT} for the definitions of $M$ and $T$.}
\begin{equation*}
(\forall\kkk)\quad
\bz_{k+1} := (1-\lambda)\bz_k + \lambda T\bz_k.
\end{equation*}
Set 
\begin{equation*}
p := \frac{1}{n-1}\big(z_{0,1}+\cdots+z_{0,n-1}\big). 
\end{equation*}
Then there exists $\bar{\bz}\in X^{n-1}$ such that 
\begin{equation*}
\bar{\bz}_k \to \bar{\bz} \in \Fix T,
\end{equation*}
and 
\begin{equation}
\label{e:210817h}
M\bz_k \to M\bar{\bz} = (P_Zp,\ldots,P_Zp) \in X^n. 
\end{equation}
In particular, 
\begin{equation}
\label{e:210817i}
P_1(Q_1\bz_k)= Q_1M\bz_k \to P_Z(p) = \tfrac{1}{n-1}P_Z\big(z_{0,1}+\cdots+z_{0,n-1}\big).
\end{equation}
Consequently, if $x_0\in X$ and 
$\bz_0 = (x_0,\ldots,x_0)\in X^{n-1}$,
then 
\begin{equation}
\label{e:yayyay}
P_1Q_1\bz_k \to P_Zx_0.
\end{equation}
\end{theorem}

\begin{proof}
Set $T_\lambda := (1-\lambda)\Id+\lambda T$ and observe that 
$(\bz_k)_\kkk = (T_\lambda^k\bz)_\kkk$. 
Hence, by \cref{c:key} and \cref{mtfixlemma},
\begin{align*}
\bz_k&\to P_{\Fix T_\lambda}\bz_0=P_{\Fix T}\bz_0\\
&=(P_Zp,\ldots,P_Zp)\oplus P_E(\bz_0), 
\end{align*}
where $E$ is as in \cref{mtfixlemma}. 
Hence, using also \cref{e:210817g}, 
\begin{align*}
Q_1M\bz_k &= P_1Q_1\bz_k \\
&\to P_1Q_1\big((P_Zp,\ldots,P_Zp)\oplus P_E(\bz_0)\big)\\
&=P_1\big(P_Zp+Q_1(P_E(\bz_0))\big)\\
&\in P_1\big(P_Zp+U_1^\perp\big)\\
&=\{P_1P_Zp\}\\
&=\{P_Zp\},
\end{align*}
i.e., $Q_1M\bz_k\to P_Zp$. 
Now \cref{e:210817a} yields 
$Q_iM\bz_k\to P_Zp$ for every $i\in\{1,\ldots,n\}$.
This yields \cref{e:210817h} and \cref{e:210817i}.

The ``Consequently'' part is clear because when 
$\bz_0$ has this special form, then $p=x_0$. 
\end{proof}

\end{document}
